
\documentclass[11pt]{article}
\usepackage[a4paper,margin=1in]{geometry}
\usepackage{amsmath,amssymb,booktabs,longtable}
\usepackage[hidelinks]{hyperref}

\title{Microneedle Superdisintegrant Film Simulator v5\\Physics, Mathematics, and Code Architecture}
\author{Generated Documentation}
\date{}

\begin{document}
\maketitle

\tableofcontents

\section{Overview}

This simulator models drug release from coated microneedles using a hybrid physics--empirical framework combining:

\begin{itemize}
\item Geometry-dependent surface area calculations
\item Capillary wicking (Lucas--Washburn theory)
\item Effective diffusion in porous TPMS structures
\item Hydration-controlled film disintegration
\item Korsmeyer--Peppas drug release kinetics
\item Burst release and sustained release decomposition
\item Excipient-dependent transport modifiers
\end{itemize}

The code is not a first-principles mechanistic model. It is a semi-empirical engineering model calibrated to trends reported in the coated microneedle, oral film, and superdisintegrant literature.

\section{Physics Model}

\subsection{Water Diffusion Through HPMC}

The simulator defines an effective water diffusivity:

\begin{equation}
D_{water}
=10^{-10}
\left(\frac{5}{\mu_{ref}}\right)^{0.28}
e^{-0.03C_{HPMC}}
\end{equation}

where:

\begin{itemize}
\item $\mu_{ref}$ = viscosity grade reference (E5, E15, E50, K4M)
\item $C_{HPMC}$ = HPMC concentration (\% w/w)
\end{itemize}

Interpretation:

\begin{itemize}
\item Higher HPMC viscosity grade lowers diffusivity.
\item Higher polymer concentration lowers diffusivity exponentially.
\item Baseline diffusivity is approximately $10^{-10}$ m$^2$/s.
\end{itemize}

\subsection{Coating Solution Viscosity}

The coating viscosity is estimated as

\begin{equation}
\mu =
\mu_{ref}
\left(
\frac{C_{HPMC}}{2}
\right)^{2.2}
\end{equation}

This is an empirical scaling relation.

\subsection{Surface Tension}

Tween-80 lowers surface tension according to

\begin{equation}
\gamma =
72 - 35 \log_{10}
\left(
1+\frac{C_{Tween}}{0.01}
\right)
\end{equation}

bounded between 28 and 72 mN/m.

\section{Geometry Module}

\subsection{Aspect Ratio}

\begin{equation}
AR = \frac{H}{R}
\end{equation}

where:

\begin{itemize}
\item $H$ = microneedle height
\item $R$ = base radius
\end{itemize}

\subsection{Tip Angle}

\begin{equation}
\theta =
\tan^{-1}\left(\frac{R}{H}\right)
\end{equation}

\subsection{TPMS Effective Diffusion}

For porous geometries:

\begin{equation}
\frac{D_{eff}}{D}
=
\frac{\varepsilon}{\tau^2}
\end{equation}

where:

\begin{itemize}
\item $\varepsilon$ = porosity
\item $\tau$ = tortuosity
\end{itemize}

This is a Millington--Quirk style correction.

\section{Surface Area Calculations}

\subsection{Planar Film}

\begin{equation}
SA = \pi R^2
\end{equation}

\subsection{Conical Microneedle}

Slant height:

\begin{equation}
s=\sqrt{H^2+(R-r_t)^2}
\end{equation}

Surface area:

\begin{equation}
SA=\pi(R+r_t)s
\end{equation}

where $r_t=5\,\mu m$ is the assumed tip radius.

\subsection{Pyramidal Microneedle}

\begin{equation}
SA = 4 \times \frac12 a s
\end{equation}

where $a=2R$.

\subsection{TPMS Structures}

The simulator multiplies the external surface area by a geometry-specific surface enhancement factor:

\begin{equation}
SA_{TPMS}=SA_{external}\times SA_{factor}
\end{equation}

to approximate internal pore surface.

\section{Drug Loading}

Coating volume:

\begin{equation}
V_{coat}=SA\times L\times f_{vol}
\end{equation}

Film density:

\begin{equation}
\rho =1.2\ g/cm^3
\end{equation}

Coating mass:

\begin{equation}
M_{coat}=V_{coat}\rho
\end{equation}

Drug mass:

\begin{equation}
M_{drug}=M_{coat}f_{drug}
\end{equation}

\section{Capillary Wicking}

The simulator implements Lucas--Washburn penetration:

\begin{equation}
x=\sqrt{Ct}
\end{equation}

with

\begin{equation}
C=
\frac{r\gamma\cos\theta}
{2\eta\tau^2}
\end{equation}

where:

\begin{itemize}
\item $r$ = effective pore radius
\item $\gamma$ = surface tension
\item $\eta$ = viscosity
\item $\theta$ = contact angle
\item $\tau$ = tortuosity
\end{itemize}

Two separate wicking calculations are performed:

\begin{enumerate}
\item Manufacturing coating uptake
\item Interstitial fluid uptake in skin
\end{enumerate}

\section{Hydration Kinetics}

Hydration rate constant:

\begin{equation}
k_{hyd}
=
\frac{D_{water}}{L^2}
\times 60
\end{equation}

modified by:

\begin{itemize}
\item Wetting factor
\item Superdisintegrant factor
\item Gel penalty
\item Surface-area-to-volume effects
\item TPMS porosity effects
\end{itemize}

Hydration fraction:

\begin{equation}
H(t)=1-e^{-k_{hyd}t}
\end{equation}

\section{Release Model}

\subsection{Korsmeyer--Peppas}

The sustained component follows

\begin{equation}
F=k_{KP}t^n
\end{equation}

where:

\begin{itemize}
\item $k_{KP}$ = release constant
\item $n$ = release exponent
\end{itemize}

\subsection{Release Mechanism Interpretation}

\begin{longtable}{ll}
\toprule
$n$ & Mechanism\\
\midrule
$n \le 0.45$ & Fickian diffusion\\
$0.45<n<0.89$ & Anomalous transport\\
$n\ge0.89$ & Case-II transport\\
\bottomrule
\end{longtable}

\subsection{Burst Release}

Burst fraction:

\begin{equation}
F_{burst}
=
B
\left(
1-e^{-k_{burst,eff}t}
\right)
\end{equation}

where

\begin{equation}
k_{burst,eff}
=
\frac{k_{hyd}k_{burst}}
{k_{hyd}+k_{burst}}
\end{equation}

(modified by gel effects in the implementation).

\subsection{Total Release}

\begin{equation}
F_{total}
=
F_{burst}
+
(1-B)F_{KP}
\end{equation}

with clipping between 0 and 100\%.

\section{Superdisintegrant Effects}

\subsection{Sodium Starch Glycolate (SSG)}

Default parameters:

\begin{itemize}
\item $k_0=0.11$
\item $n_0=0.65$
\item Burst = 22\%
\item Gel threshold = 8\%
\end{itemize}

\subsection{Croscarmellose Sodium (CCS)}

\begin{itemize}
\item $k_0=0.14$
\item $n_0=0.55$
\item Burst = 27\%
\item Gel threshold = 5\%
\end{itemize}

\subsection{PVPP}

\begin{itemize}
\item $k_0=0.23$
\item $n_0=0.40$
\item Burst = 38\%
\item Essentially no gel barrier
\end{itemize}

\subsection{Gel Penalty}

Above threshold:

\begin{equation}
Penalty = 0.75^{(C-C_{crit})}
\end{equation}

This reduces hydration and release.

\section{TPMS Physics}

The simulator contains:

\begin{itemize}
\item Gyroid
\item Primitive
\item Diamond
\item Cone + TPMS core
\end{itemize}

Implicit equations used for visualization:

\subsection{Gyroid}

\begin{equation}
\sin x\cos y+
\sin y\cos z+
\sin z\cos x=0
\end{equation}

\subsection{Primitive}

\begin{equation}
\cos x+\cos y+\cos z=0
\end{equation}

\subsection{Diamond}

\begin{equation}
\sin x\sin y\sin z
+\sin x\cos y\cos z
+\cos x\sin y\cos z
+\cos x\cos y\sin z=0
\end{equation}

\section{Code Architecture}

\begin{enumerate}
\item User enters formulation and geometry.
\item \texttt{materialModel()} computes viscosity, surface tension, diffusivity, and kinetic modifiers.
\item \texttt{geometryModel()} computes area, volume, drug loading, TPMS corrections.
\item \texttt{simulateWicking()} computes capillary penetration versus time.
\item \texttt{simulateRelease()} computes hydration, burst release, sustained release, and total release.
\item Results are rendered using Recharts.
\item TPMS structures are visualized with Three.js.
\end{enumerate}

\section{Important Assumptions}

\begin{itemize}
\item Uniform coating density = 1.2 g/cm$^3$
\item Constant diffusivity
\item No skin metabolism
\item No drug degradation
\item No concentration-dependent diffusivity
\item Porous structures represented using effective-medium theory
\item Korsmeyer--Peppas parameters are empirical calibration constants
\end{itemize}

\section{Likely Literature Sources}

The code appears to draw concepts from:

\begin{enumerate}
\item Korsmeyer RW et al. Mechanisms of solute release from porous hydrophilic polymers.
\item Peppas NA and Sahlin JJ. Mathematical modeling of drug release.
\item Lucas and Washburn capillary penetration theory.
\item Millington and Quirk tortuosity correction.
\item Microneedle coating and dissolving microneedle literature.
\item HPMC film-forming and oral film formulation studies.
\item CCS, SSG, and PVPP superdisintegrant characterization studies.
\end{enumerate}

\end{document}
